O Semiárido brasileiro apresenta elevada variabilidade das precipitações, perdas expressivas por evaporação e secas recorrentes, condições que tornam a operação de reservatórios relevante para a segurança hídrica. Este trabalho teve como objetivo desenvolver e avaliar um sistema web de suporte à decisão baseado na simulação mensal do balanço hídrico para a análise da operação de reservatórios individuais e de hidrossistemas no Ceará. O simulador considera afluência, evaporação pela área média do espelho d'água, retirada, vertimento e transferências; representa as operações Individual, Série e Paralelo; e aplica curvas-guia mensais com racionamento. A interface, desenvolvida em React e integrada a uma interface de programação de aplicações em FastAPI e a uma base SQLite com 155 reservatórios, permite configurar as análises, visualizar indicadores e séries mensais e exportar os resultados. A implementação foi verificada por comparação com um cálculo independente em planilha, por conferências das regras de racionamento e por nove casos controlados, abrangendo os reservatórios principais de 25 hidrossistemas e 33.108 balanços mensais. Planilha e simulador coincidiram em todos os hidrossistemas, com diferenças nulas dentro da tolerância numérica de 10\textsuperscript{-9}~hm\textsuperscript{3}, e todos os casos controlados foram aprovados. Na análise de sensibilidade, a demanda foi o parâmetro de maior influência sobre as falhas, a evaporação teve efeito moderado e o volume inicial, efeito desprezível sobre a série longa; o gatilho e a vazão de transferência apresentaram limiares de segurança para o reservatório receptor. As aplicações mostraram que as curvas-guia eliminaram as 142 falhas que ocorreriam em Mundaú sem racionamento, à custa de 332 meses com retirada reduzida; que as falhas em Carnaubal--Barragem do Batalhão ocorreram quando a unidade secundária não tinha volume para assumir a demanda conjunta; e que, em Fogareiro--Quixeramobim, a regra de gatilho produziu alternância da transferência em meses consecutivos. Os testes realizados demonstraram consistência entre as equações implementadas e as saídas produzidas nos casos controlados; como não houve comparação com volumes observados, os resultados não constituem validação da representação da operação real.

% Separe as palavras-chave por ponto
\palavraschave{recursos hídricos; operação de reservatórios; sistema de suporte à decisão; simulação hidrológica; verificação de modelos.}
